\section{Supplementary assay and fitting details}
\label{app:fit-details}

The models and fitting procedures use the \texttt{flipv3} module of
SciStanPy at revision
\href{https://github.com/microsoft/SciStanPy/tree/09c439f675ead70f851d814242f153ecb7712113/flipv3}{\texttt{09c439f675ea}},
the version pinned by the companion analysis repository. We first describe
the observations and model equations, then report the fitting records and
available diagnostics. Source defaults, requested command settings, and
values confirmed by saved artifacts are distinguished throughout. TODO
entries mark missing records or checks.

\subsection{Experimental records and candidate identities}
\label{app:experimental-records}

The source studies are the PDZ interaction assays of \citet{zarin2024}, the
TrpB active-site libraries of \citet{johnston2024}, and the NucB engineering
generations of \citet{thomas2025}. We retain their landscape labels rather
than counting fitted models as new experiments. The PDZ inputs
are cis, trans-1, trans-2, CRIPT-C, and CRIPT-N. TrpB contains three-site
libraries A--I and a four-site library; the labels refer to distinct variant
libraries, not different resampling seeds. NucB generations G1--G4 are
successive engineering collections with different screening protocols.

\paragraph{Physical assays and reported labels.}
The PDZ experiment measures binding between PSD95-PDZ3 and CRIPT through
protein complementation in yeast. Its source protocol
(\href{https://www.ncbi.nlm.nih.gov/geo/query/acc.cgi?acc=GSE265816}{GSE265816})
documents three biological growth competitions for each library. Input
populations are grown without the selective drug; the output competition
uses methotrexate at $200\,\mu\mathrm{g}/\mathrm{mL}$ for approximately five
generations. DNA from the input and output cultures is sequenced, and DiMSum
produces replicate-combined binding-dependent fitness estimates and standard
errors. The deposited protocol also describes normalization between library
blocks using stop-codon and synonymous variants, followed by linear scaling
between libraries using shared variants. These source score-processing steps
are distinct from fitting our count models.

The TrpB libraries use an engineered background derived from the
\emph{Thermotoga maritima} enzyme. Bacteria lacking the native
\textit{trpA} and \textit{trpB} genes express a library of introduced TrpB
sequences. With indole supplied, synthesis of tryptophan supports growth;
sequencing observes the resulting changes in library composition
\citep{johnston2024}. The nine three-site libraries together cover twenty
positions, while the four-site library varies positions 183, 184, 227, and
228. A stop codon terminates translation prematurely, providing a
nonfunctional control for the source growth-score normalization. The count
model does not estimate separate catalytic, expression, and stability effects.

The NucB generations are successive rounds of engineering, not four biological
replicates of a fixed library. Their input and selected populations have
different fluorescence thresholds and sequencing structures
\citep{thomas2025}. The source release contains pre-sort and post-sort counts,
enrichment estimates relative to known controls, and activity labels combined
across rounds. Our model uses the generation-specific count observations;
it does not treat the combined categorical activity labels as count data.
Source low-throughput measurements of purified enzymes are separate validation
experiments and do not enter the rank calculations here.

\note{TODO: Complete the residue-level TrpB three-site designs. Specify the
constructs and mutated sites in PDZ cis, trans-1, trans-2, CRIPT-C, and CRIPT-N.
Verify the NucB fluorescent substrate, compartment protocol, and incubation
conditions against the full experimental methods. Add the per-task
correspondence to FLIP2 without treating its subsets as additional experiments.}

\paragraph{Count axes.}
The model input arrays preserve the following observation structures. A
repeated sequencing library is distinguished from a new biological population.
The last array dimension is always the fitted variant axis; source controls
can belong to this axis without being eligible ranking candidates.
\begin{center}
\small
\begin{tabular}{lll}
\toprule
Assay & Input observations & Output observations \\
\midrule
PDZ & Three replicate vectors & Three paired vectors \\
TrpB A--C & One shared vector & One vector per timepoint \\
TrpB D--I & One shared vector & Two replicates per timepoint \\
TrpB four-site & Two replicate vectors & Two replicates per timepoint \\
NucB G1 & Five vectors, three populations & Seven selected count vectors \\
NucB G2 & Three vectors, two populations & Three selected count vectors \\
NucB G3--G4 & One vector per generation & Four selected count vectors \\
\bottomrule
\end{tabular}
\end{center}
Generation 1 sequences the first two input populations twice and the third
once; generation 2 sequences the first input twice and the second once.

\subsection{PDZ count model and latent fitness}
\label{app:pdz-model}

\paragraph{Observation model.}
Let $M$ denote the number of variants on the fitted count axis, including
any rows that are excluded from ranking. This is distinct from the number
$N$ of ranking-eligible variants in the main text. For each of the three
PDZ replicates $r\in\{1,2,3\}$, let
$\mathbf n_r^{\mathrm{in}}=(n_{1r}^{\mathrm{in}},\ldots,n_{Mr}^{\mathrm{in}})$
and $\mathbf n_r^{\mathrm{out}}$ be its input and output count vectors. Write
$D_r^{\mathrm{in}}:=\sum_{i=1}^M n_{ir}^{\mathrm{in}}$ and
$D_r^{\mathrm{out}}:=\sum_{i=1}^M n_{ir}^{\mathrm{out}}$ for their observed
sequencing totals. The model conditions on these totals.

Let $\boldsymbol\pi_r^0$ be the vector of input probabilities for replicate
$r$, with nonnegative entries summing to one. A growth function gives positive
unnormalized output abundances $a_{ir}$ from the input probabilities and
replicate-level growth parameters. Define
$\pi_{ir}^1:=a_{ir}/\sum_{j=1}^M a_{jr}$. The count model is
\begin{align}
\mathbf n_r^{\mathrm{in}}\mid\theta
&\sim\operatorname{Multinomial}(D_r^{\mathrm{in}},\boldsymbol\pi_r^0),\\
\mathbf n_r^{\mathrm{out}}\mid\theta
&\sim\operatorname{Multinomial}(D_r^{\mathrm{out}},\boldsymbol\pi_r^1).
\end{align}
Here $\theta$ includes the input probabilities, the growth parameters, and
any shared parameters. Conditional on those parameters, the input and output
observations contribute separate likelihood terms. Within a library, the
multinomial distribution retains dependence between variant counts through
their fixed total. The implementation works with log probabilities and
normalizes the output abundances on that scale.

The source defines separate input probability vectors for the three PDZ
replicates and a hierarchy linking their variant-specific growth rates.
Each input vector has a symmetric Dirichlet prior,
$\boldsymbol\pi_r^0\sim\operatorname{Dirichlet}(\alpha_0\mathbf1_M)$,
where $\mathbf1_M$ is a vector of $M$ ones and $\alpha_0>0$ is the
per-variant concentration. The source class \texttt{ExpDirichlet} stores
$\log\boldsymbol\pi_r^0$; its name denotes a log parameterization, not a
different input-count likelihood.

\paragraph{Growth functions.}
Let $b_{ir}$ be the effective growth parameter for variant $i$ in replicate
$r$, and let $\tau\geq0$ be elapsed time on the model's scale. PDZ has one
post-selection observation, for which $\tau=1$. For exponential growth the
unnormalized abundance is
\[
a_{ir}(\tau):=\pi_{ir}^0\exp(b_{ir}\tau).
\]
For sigmoid growth, the implementation instead uses
\[
a_{ir}(\tau):=\pi_{ir}^0
\frac{1+\exp(b_{ir}c)}{1+\exp[b_{ir}(c-\tau)]},
\qquad c\sim\operatorname{Gamma}(a_c,b_c).
\]
Here $c$ is a shared positive offset. We use the shape--rate convention
for every Gamma distribution in this supplement, so its mean is shape
divided by rate. Both growth functions start at $a_{ir}(0)=\pi_{ir}^0$.
Output probabilities normalize $a_{ir}(\tau)$ over variants separately
for each observed library. We use these functions as assay models of abundance
change, not as mechanistic descriptions of every stage of the biological assay.

\paragraph{Exponential and Lomax rate families.}
For each variant, let $m_i>0$ be its central growth parameter and
$\eta_i:=\log m_i$. The Exponential specification assigns
$m_i\sim\operatorname{Exponential}(\beta)$, with rate $\beta$ and mean
$1/\beta$. The Lomax specification assigns a distribution with scale
$\lambda>0$, shape $a_L>0$, and density
\[
p(m_i\mid\lambda,a_L)=\frac{a_L}{\lambda}
(1+m_i/\lambda)^{-(a_L+1)},\qquad m_i>0.
\]
SciStanPy's \texttt{ExpExponential} and \texttt{ExpLomax} variables store
$\eta_i$, including the appropriate transformed density. They do not mean
that $\eta_i$ itself has an ordinary Exponential or Lomax distribution.

There are two available replicate models. The normal option uses
$b_{ir}\mid m_i\sim\operatorname{Normal}(m_i,\sigma_b^2)$. Its central
parameter is the conditional mean rate, and the replicate rates may be
negative. The lognormal option uses
$\log b_{ir}\mid\eta_i\sim\operatorname{Normal}(\eta_i,\sigma_b^2)$.
Here $m_i$ is the conditional median rate; the conditional mean is
$m_i\exp(\sigma_b^2/2)$. Since $\sigma_b$ is shared, both induce the same
ordering of these central rates. The hierarchy can be represented using
standard-normal latent variables or using the replicate variables directly.
The command flag \texttt{--centered} selects the latter representation where
the model exposes that option. The command records, rather than the default
alone, determine whether a run uses normal or lognormal replicate variation.

\paragraph{Gamma inverse-rate family.}
The Gamma specification assigns an inverse-rate parameter
$\kappa_i\sim\operatorname{Gamma}(a_\kappa,b_\kappa)$ and
$b_{ir}\mid\kappa_i\sim\operatorname{Exponential}(\kappa_i)$.
The conditional mean growth rate is $m_i:=1/\kappa_i$. The stored
parameter is \texttt{inv\_r\_mean}, and the score extractor applies
$h_i(\theta):=-\log\kappa_i$. For the Exponential and Lomax hierarchies,
the stored variable is \texttt{log\_r\_mean} and
$h_i(\theta):=\eta_i$. Thus the score labeled \texttt{log\_mean\_growth}
orders the central growth parameter; in the lognormal hierarchy it is not
literally the logarithm of the arithmetic mean replicate rate.

\paragraph{Hyperparameters.}
The source defaults are $\alpha_0=0.1$, $\beta=10$, $\lambda=1$,
$a_L=2$, $\sigma_b=0.1$, $(a_\kappa,b_\kappa)=(2,0.5)$, and
$(a_c,b_c)=(4,8)$. The job commands can override them. In particular,
the recorded PDZ commands set $\alpha_0=5$ and select normal replicate
variation for the Exponential and Lomax hierarchies. We retain per-run
command settings in the fit inventory rather than replacing overrides with
the class defaults.

\subsection{TrpB count model and latent fitness}
\label{app:trpb-model}

\paragraph{Time and observation structure.}
TrpB uses the same growth functions and multinomial observation model as
PDZ, but has multiple post-selection timepoints. Let $t_\ell$ be an
observed time, indexed by $\ell$, and define
$\tau_\ell:=t_\ell/\max_{\ell'}t_{\ell'}$. Effective growth rates are
therefore expressed per normalized assay duration, not per hour. The model
factory removes optical-density columns before model construction; the
likelihood uses sequencing counts and normalized times.

Libraries A--C have one input count vector and one output vector per timepoint.
They use a single input probability vector $\boldsymbol\pi^0$ and one rate
$b_i$ per variant across the observed times. The prior is Exponential or
Lomax on $b_i$, with the definitions given above; the flat Gamma alternative
is not implemented. Their latent score is $h_i(\theta):=\log b_i$, obtained
from the stored positive variable \texttt{r}.

Libraries D--I have one shared input count vector and replicate-specific
output vectors at each retained time. They use one
$\boldsymbol\pi^0$ across output replicates and hierarchical $b_{ir}$
rates, with the Exponential, Lomax, or Gamma inverse-rate construction
above. For each time and replicate, the output vector has a multinomial
likelihood with its own observed total and normalized grown probabilities.
The single input observation contributes one likelihood term; it is not
repeated once per output replicate. Their score is $\eta_i$ or
$-\log\kappa_i$, as for PDZ.

The four-site model has two input count vectors and replicate-specific
outputs. Both input vectors share one latent input probability vector, while
the output rates use the same hierarchical families. The prepared paper
collection contains optimized-fit count resampling, but no parameter-posterior
or posterior-count analysis, for this landscape.

\paragraph{Model and scoring timepoints.}
The model uses the times supplied by the data loader; the count-derived
fitness calculation can use a different retained subset. In particular,
the four-site loader supplies 12, 16, 20, 24, 28, and 36 hours, while the
deposited count-score reconstruction excludes the 12-hour observation.
Appendix~\ref{app:scores} distinguishes those score rules from the count
likelihood. The input Dirichlet concentration and other hyperparameters have
the defaults above, with overrides and centered/noncentered choices listed
for each job.

\subsection{NucB gate model and latent fitness}
\label{app:nuc-model}

\paragraph{Latent fluorescence.}
For variant $i$, let $L_i$ be a log-fluorescence location and
$v_i:=\exp(L_i)>0$ its fluorescence-scale center. The Exponential family
assigns $v_i\sim\operatorname{Exponential}(\beta)$; the Lomax family
assigns the scale--shape distribution defined above. Both use a log
parameterization internally. The paper's parameter score is
$h_i(\theta):=L_i$, not a count enrichment and not raw measured fluorescence.
For generation 1 the stored variable is
\texttt{absolute\_mean\_log\_fluorescence}; for generations 2--4 it is
\texttt{log\_fluorescence}. The extractor retains those values without
exponentiating them. The data-product label \texttt{model\_fluorescence}
therefore denotes a log-scale parameter in these analyses.

\paragraph{Gate probabilities.}
Let $z_g>0$ be the fluorescence threshold for gate $g$, and
$\sigma_F>0$ a shared spread parameter. For a lognormal observation model,
the probability of exceeding the threshold for a variant with location $L$
is
\[
s_g(L,\sigma_F):=1-\Phi\!\left(\frac{\log z_g-L}{\sigma_F}\right),
\]
where $\Phi$ is the standard-normal cumulative distribution function.
Generations 3 and 4 also expose a normal option,
$s_g(L,\sigma_F):=1-\Phi[(z_g-\exp L)/\sigma_F]$.
The recorded generation-3 command uses that normal option; generations 1
and 2 use the lognormal survival model, and the recorded generation-4
command leaves the lognormal default in place.

For gate $g$ arising from input population $r$, the variant probabilities
are
\[
\pi_{ig}^{\mathrm{gate}}:=
\frac{\pi_{ir}^0s_g(L_{ir},\sigma_F)}
{\sum_{j=1}^M\pi_{jr}^0s_g(L_{jr},\sigma_F)}.
\]
Here $L_{ir}$ is the population-specific location when the model includes
one, and otherwise equals $L_i$. Each observed gate count vector is
multinomial with these probabilities and its observed sequencing total.
The gate likelihoods use survival probabilities above thresholds, not a
partition of all cells into mutually exclusive fluorescence bins. Repeated
sequencing of an input or selected population shares its underlying probability
vector but contributes its own count likelihood and total.

\paragraph{Generations.}
Generation 1 has three input-population vectors,
$\boldsymbol\pi_r^0\sim\operatorname{Dirichlet}(\alpha_0\mathbf1_M)$,
and population-specific log locations
$L_{ir}\mid L_i,\sigma_E\sim\operatorname{Normal}(L_i,\sigma_E^2)$.
The experimental spread satisfies
$\sigma_E\sim\operatorname{HalfNormal}(s_E)$, meaning the absolute value
of a mean-zero normal draw with standard deviation $s_E$. Low- and high-threshold
selected counts use the corresponding location and input population.
Generation 2 has two input vectors and a shared $L_i$ across gates; the
86 and 97.5 gates use the first input population, and the 93 gate uses the
second. Generations 3 and 4 each have one input vector and four thresholded
output count vectors, all sharing $L_i$.

All generations assign
$\sigma_F\sim\operatorname{Gamma}(a_F,b_F)$. The source names this
parameter \texttt{codon\_noise}; that name is a modeling interpretation,
not evidence that the observed spread can be attributed only to codon usage.
Defaults are $(a_F,b_F)=(2,2)$ and $s_E=0.5$. The command records use
$\alpha_0=5$ for NucB and request Exponential rates $10,12.5,15,17.5$ in
generations 1--4, respectively, through the flag \texttt{exp\_beta}.
However, the inspected factory forwards that keyword while the Exponential
mixin accepts \texttt{beta}; the unmatched keyword is absorbed without
overriding the default. Under this source revision the effective rate is
therefore 10, unless \texttt{beta} is passed through another entry point.
We record the requested setting and this source-level behavior separately.
The Lomax shape keyword matches its constructor argument; the recorded shapes
are $2,4,6,8$, respectively, at scale one.
Only specifications present in the curated fit inventory contribute to the
paper's results.

\paragraph{Threshold values.}
The immutable data loader supplies the following thresholds on its fluorescence
scale. These are assay calibration constants, not fitted rank cutoffs.
For generation 1, the three pairs correspond to the three input populations.
For generations 2--4, values follow the gate order used in the count arrays.
\begin{center}
\small
\begin{tabular}{lll}
\toprule
Generation & Gate labels & Threshold values \\
\midrule
G1 & Low, by population & 0.100, 0.115, 0.100 \\
G1 & High, by population & 0.501, 0.600, 0.370 \\
G2 & 86, 93, 97.5 & 0.201, 0.318, 0.385 \\
G3 & 59, 80, 95, 99 & 0.118, 0.217, 0.285, 0.346 \\
G4 & 70, 90, 98, 99.5 & 0.124, 0.199, 0.306, 0.384 \\
\bottomrule
\end{tabular}
\end{center}

\paragraph{Controls and interpretation.}
Appended negative-control rows enter the multinomial normalization and fitted
axis but not the candidate ranking. Gate thresholds determine the connection
between the latent location and the assay; none of these parameters is a
universal biochemical fitness measure. The highest-gate enrichment score used
for the count-resampling analyses is defined separately in
Appendix~\ref{app:scores}. It is not identical to $L_i$, even though both
are oriented so that higher scores correspond to higher inferred assay activity.

\subsection{Optimization and posterior simulation}
\label{app:inference-details}

\paragraph{Posterior simulation.}
The \texttt{mcmc\_flip.py} entry point constructs the selected model,
generates and compiles its Stan program, and runs Hamiltonian Monte Carlo
through CmdStan. It then converts samples and generated quantities to NetCDF
and runs the diagnostic routines. The command-line defaults are eight chains,
2,000 warmup iterations and 2,000 retained iterations per chain, with seed
1025. The job command list overrides warmup and retained iterations for some
models. The saved artifact attributes, rather than the defaults, are the
source for actual chain counts, iteration counts, and Stan versions in the
fit inventory. The representative artifacts inspected here record Stan 2.38.0
and maximum tree depth 10. Absence of an explicit adaptation option in the
launcher is not a record of its effective per-fit value.

\paragraph{Optimized fit.}
The source calls the second entry point \texttt{mle\_flip.py}. Its
PyTorch objective sums log densities of both
observed nodes and latent parameter nodes. In the parameterization evaluated
by that backend, it maximizes a joint-density objective of the form
\[
\hat\theta:=\arg\max_\theta
\left\{\sum_{o\in\mathcal O}\log p(d_o\mid\operatorname{pa}(o))
+\sum_{v\in\mathcal V}\log p(\theta_v\mid\operatorname{pa}(v))\right\}.
\]
Here $\mathcal O$ indexes observed count nodes with values $d_o$,
$\mathcal V$ indexes latent stochastic nodes, and
$\operatorname{pa}(\cdot)$ denotes a node's parents in the model. Thus the
fitted mode includes priors and hierarchical densities; it is not the
maximizer of the count likelihood alone. Density modes also depend on the
chosen parameterization. We use \emph{optimized-fit count resampling} in the
prose and preserve \texttt{mle\_predictive} as the existing data identifier.
This correction changes the description, not the stored fits or results.

Optimization uses Adam, with source defaults of learning rate $0.001$, at
most 100,000 full-data steps, and stopping after 10 steps without an improvement
in the recorded objective. The routine retains a saved best state when early
stopping is triggered and writes the loss trajectory. The launcher passes
seed 1025 unless overridden. These are reproducible source settings; the
inventory flags where a historical artifact or loss record is unavailable
to confirm an individual execution. Optimization convergence requires the
loss and restart records listed below.

\paragraph{Resampled counts.}
Posterior-count samples are generated using sampled model parameters and
the corresponding fixed library totals. The optimized-fit workflow exports
the fitted distributions and produces count realizations using its sampling
routine without reoptimizing on each realization. Appendix~\ref{app:scores}
then maps those counts to the same family-specific reconstruction functions
used throughout the rank analysis. Preserving one joint draw across all
variant and observation axes is necessary for coherent rank uncertainty.

\subsection{Convergence and predictive checks}
\label{app:fit-checks}

\paragraph{Available sampler evidence.}
All 49 posterior artifacts were accessible for this inventory. Their recorded
chain count, warmup, retained iterations, and seed agree with the corresponding
job commands. All have eight chains, seed 1025, Stan 2.38.0, and a recorded
maximum tree depth of 10. None records a divergent transition. Six fits record
tree-depth saturation: TrpB jobs 20, 25, 26, 38, 43, and 44 have 4,000,
16,000, 16,000, 1,692, 12,454, and 14,000 saturated transitions, respectively.
These are warnings about trajectory length and computational efficiency;
zero divergences does not remove the need to examine them.

Forty artifacts include convergence and effective sample size summaries for
the stored score parameter. Across their fitted variant axes, the largest
finite recorded convergence statistic $\hat R$ is approximately 1.0067,
the smallest bulk effective sample size is 1,876.9,
and the smallest tail effective sample size is 889.2. These diagnostics apply
to the native stored score parameter, including fitted controls when present.
They are not effective sample sizes for ranks or top-$k$ membership. The nine
raw-artifact fallbacks lack these stored summaries and are shown as unavailable.
The code's variable-level diagnostic
defaults flag $\hat R\geq1.01$ and effective sample sizes at most 100 per
chain. We report the recorded values rather than assigning one pass/fail
label to the whole analysis.

\paragraph{Energy diagnostic.}
For each saved chain with energy sequence $E_1,\ldots,E_T$, we compute the
ratio
\[
B:=\frac{(T-1)^{-1}\sum_{t=2}^T(E_t-E_{t-1})^2}
{T^{-1}\sum_{t=1}^T(E_t-\overline E)^2},
\qquad \overline E:=T^{-1}\sum_{t=1}^T E_t.
\]
This is the energy-based fraction of missing information (BFMI) estimate
using the displayed variance convention. The minimum over all inspected
chains is 0.637. We calculate this from energy traces because the inspected
source's stored flag named \texttt{low\_ebfmi} compares individual energy
values with a threshold rather than calculating this chain-level ratio.
The inventory does not interpret that flag as a valid BFMI test.

\paragraph{What these checks do not establish.}
The inventory records diagnostic availability and values for the retained
fits without applying new exclusions. Sampler diagnostics assess simulation
from the fitted distribution, not whether that distribution describes the
experiment. Predictive checks remain necessary, particularly for specifications
that produce very different rank uncertainty from the same counts.

\subsubsection{Reserved predictive-check exhibits}
\note{TODO: Add per-fit observed-versus-predictive checks for count support,
zero counts, replicate or timepoint variation, and gate enrichment. Define
each discrepancy statistic and show the relevant biological examples. This
subsection is intentionally open; the sampler inventory is not a substitute
for these checks.}

\subsubsection{Reserved exception and optimization checks}
\note{TODO: Recompute missing native-score diagnostics for raw fallbacks;
quantify rank- and inclusion-specific Monte Carlo precision; review tree-depth
saturation and the strongly specification-sensitive PDZ fits; recover the
optimization loss curves and any restart evidence for all optimized fits.
Document the original retention decisions separately from any new sensitivity
analysis. No new fits or filters have been applied in this draft.}

\subsection{Score reconstruction and missing-data audit}
\label{app:score-audit}
The TrpB and NucB reconstructions reproduce the source finite scores and
coverage at the numerical tolerances stated in Appendix~\ref{app:scores}.
The fit mapping records retained samples and exclusions. The PDZ reconstruction
is an approximation because the exact original DiMSum configuration has
not been recovered.
\note{TODO: Add the per-assay reconstruction comparisons, original DiMSum
revision and configuration if recovered, finite-score support distributions,
and sensitivity of reported-rank summaries to missingness and ties.}

\subsection{Reproducibility and complete supplementary exhibits}
\label{app:reproducibility}

The two tables below enumerate all 49 posterior fits and all 60 optimized
fits used by the 158 analysis runs. Each posterior fit supplies both a
parameter and a posterior-count view. The separate machine-readable inventory
records each analysis identifier, its fitted artifact, its source job command,
all parsed options, retained samples and exclusions, and available native-score
diagnostics. Source hashes identify the model and launcher files used for this
description. Historical job commands are records of requested settings; saved
artifact attributes confirm the reported sampling counts, but do not embed
every prior option or the complete software environment.

Model codes combine the rate-family initial (E: Exponential, L: Lomax,
G: Gamma inverse-rate) and growth-function initial (E: exponential,
S: sigmoid). NucB has only the fluorescence-family code E or L. A source
artifact may have been retired after its prepared score matrix was validated;
this is distinct from a missing paper analysis. Three original optimized
artifacts remain locally available, while all 60 prepared optimized-count
score matrices are available for the rank analysis.

\begingroup
\scriptsize
\setlength{\tabcolsep}{3pt}
\begin{longtable}[c]{lllllllllll}
\caption{All posterior fits used by the parameter and posterior-count analyses. Iterations are per chain from artifact metadata. D denotes a diagnosed artifact; R a raw artifact. The maximum stored native-score R-hat uses all fitted variants, not only the decision cohort; -- means unavailable. Depth counts samples reaching the recorded maximum tree depth. BFMI is recomputed from energy traces.}\label{tab:hmc-fits}\\
\toprule
Landscape & Model & Job & Chains & Warmup & Draws & File & Diverg. & Depth & $\max\hat R$ & BFMI \\
\midrule
\endfirsthead
\toprule
Landscape & Model & Job & Chains & Warmup & Draws & File & Diverg. & Depth & $\max\hat R$ & BFMI \\
\midrule
\endhead
\bottomrule
\endfoot
nuc-G1 & E & 55 & 8 & 4000 & 4000 & D & 0 & 0 & 1.0026 & 0.708 \\
nuc-G1 & L & 56 & 8 & 2000 & 4000 & R & 0 & 0 & -- & 0.739 \\
nuc-G2 & E & 57 & 8 & 2000 & 2000 & D & 0 & 0 & 1.0054 & 0.883 \\
nuc-G2 & L & 58 & 8 & 4000 & 3000 & D & 0 & 0 & 1.0037 & 0.901 \\
nuc-G3 & E & 59 & 8 & 2000 & 2000 & D & 0 & 0 & 1.0067 & 0.895 \\
nuc-G4 & E & 61 & 8 & 2000 & 2000 & D & 0 & 0 & 1.0028 & 0.885 \\
nuc-G4 & L & 62 & 8 & 2000 & 2000 & D & 0 & 0 & 1.0046 & 0.89 \\
pdz-cis & EE & 1 & 8 & 6000 & 2000 & D & 0 & 0 & 1.0024 & 0.939 \\
pdz-cis & ES & 4 & 8 & 6000 & 2000 & D & 0 & 0 & 1.0027 & 0.937 \\
pdz-cis & LE & 3 & 8 & 6000 & 2000 & D & 0 & 0 & 1.0028 & 0.893 \\
pdz-cript-c & EE & 70 & 8 & 2000 & 2000 & D & 0 & 0 & 1.0040 & 0.949 \\
pdz-cript-c & ES & 73 & 8 & 2000 & 3000 & D & 0 & 0 & 1.0026 & 0.924 \\
pdz-cript-c & GE & 72 & 8 & 2000 & 3000 & R & 0 & 0 & -- & 0.847 \\
pdz-cript-c & LE & 71 & 8 & 2000 & 3000 & R & 0 & 0 & -- & 0.908 \\
pdz-cript-n & GE & 78 & 8 & 4000 & 3000 & R & 0 & 0 & -- & 0.824 \\
pdz-cript-n & GS & 81 & 8 & 2000 & 2000 & D & 0 & 0 & 1.0027 & 0.86 \\
pdz-cript-n & LS & 80 & 8 & 2000 & 2000 & D & 0 & 0 & 1.0039 & 0.899 \\
pdz-trans-1 & EE & 82 & 8 & 2000 & 2000 & R & 0 & 0 & -- & 0.926 \\
pdz-trans-1 & GE & 84 & 8 & 2000 & 2000 & D & 0 & 0 & 1.0035 & 0.852 \\
pdz-trans-2 & EE & 88 & 8 & 2000 & 2000 & D & 0 & 0 & 1.0035 & 0.929 \\
pdz-trans-2 & GS & 93 & 8 & 2000 & 2000 & D & 0 & 0 & 1.0029 & 0.867 \\
trpb-libA & EE & 7 & 8 & 2000 & 2000 & D & 0 & 0 & 1.0027 & 0.787 \\
trpb-libA & ES & 9 & 8 & 2000 & 2000 & D & 0 & 0 & 1.0027 & 0.821 \\
trpb-libA & LE & 8 & 8 & 2000 & 2000 & D & 0 & 0 & 1.0029 & 0.644 \\
trpb-libB & EE & 11 & 8 & 2000 & 2000 & D & 0 & 0 & 1.0023 & 0.811 \\
trpb-libB & ES & 13 & 8 & 2000 & 2000 & D & 0 & 0 & 1.0028 & 0.758 \\
trpb-libB & LE & 12 & 8 & 2000 & 2000 & D & 0 & 0 & 1.0026 & 0.673 \\
trpb-libC & EE & 15 & 8 & 2000 & 2000 & R & 0 & 0 & -- & 0.775 \\
trpb-libC & ES & 17 & 8 & 2000 & 2000 & D & 0 & 0 & 1.0024 & 0.805 \\
trpb-libC & LE & 16 & 8 & 2000 & 2000 & D & 0 & 0 & 1.0027 & 0.637 \\
trpb-libD & EE & 19 & 8 & 2000 & 2000 & D & 0 & 0 & 1.0024 & 0.868 \\
trpb-libD & GE & 21 & 8 & 2000 & 2000 & D & 0 & 0 & 1.0022 & 0.708 \\
trpb-libD & LE & 20 & 8 & 2000 & 2000 & D & 0 & 4000 & 1.0041 & 0.815 \\
trpb-libE & EE & 25 & 8 & 6000 & 2000 & R & 0 & 16000 & -- & 0.812 \\
trpb-libE & GE & 27 & 8 & 6000 & 2000 & D & 0 & 0 & 1.0025 & 0.72 \\
trpb-libE & LE & 26 & 8 & 8000 & 2000 & D & 0 & 16000 & 1.0025 & 0.802 \\
trpb-libF & EE & 31 & 8 & 2000 & 2000 & D & 0 & 0 & 1.0020 & 0.908 \\
trpb-libF & ES & 34 & 8 & 2000 & 2000 & D & 0 & 0 & 1.0022 & 0.846 \\
trpb-libF & GE & 33 & 8 & 2000 & 2000 & D & 0 & 0 & 1.0018 & 0.701 \\
trpb-libF & LE & 32 & 8 & 2000 & 2000 & D & 0 & 0 & 1.0026 & 0.808 \\
trpb-libG & EE & 37 & 8 & 6000 & 2000 & R & 0 & 0 & -- & 0.852 \\
trpb-libG & GE & 39 & 8 & 4000 & 2000 & D & 0 & 0 & 1.0025 & 0.717 \\
trpb-libG & LE & 38 & 8 & 4000 & 2000 & D & 0 & 1692 & 1.0028 & 0.775 \\
trpb-libH & EE & 43 & 8 & 2000 & 2000 & R & 0 & 12454 & -- & 0.88 \\
trpb-libH & GE & 45 & 8 & 4000 & 2000 & D & 0 & 0 & 1.0023 & 0.71 \\
trpb-libH & LE & 44 & 8 & 2000 & 2000 & D & 0 & 14000 & 1.0028 & 0.8 \\
trpb-libI & EE & 49 & 8 & 2000 & 2000 & D & 0 & 0 & 1.0025 & 0.885 \\
trpb-libI & GE & 51 & 8 & 2000 & 2000 & D & 0 & 0 & 1.0025 & 0.767 \\
trpb-libI & LE & 50 & 8 & 2000 & 2000 & D & 0 & 0 & 1.0027 & 0.84 \\
\end{longtable}
\endgroup
\begingroup
\scriptsize
\setlength{\tabcolsep}{3pt}
\begin{longtable}[c]{lllllll}
\caption{All optimized fits used by the count-resampling analyses. Concentration, observation-family option, and centered (C) or noncentered (NC) flag come from the recorded job commands; flags may not affect every model class. Available indicates whether the original fitted artifact remains locally accessible, not whether its prepared scores are available. Model codes are defined in the text.}\label{tab:optimized-fits}\\
\toprule
Landscape & Model & Job & $\alpha_0$ & Variation option & Flag & Available \\
\midrule
\endfirsthead
\toprule
Landscape & Model & Job & $\alpha_0$ & Variation option & Flag & Available \\
\midrule
\endhead
\bottomrule
\endfoot
nuc-G1 & E & 55 & 5 & lognormal & C & no \\
nuc-G1 & L & 56 & 5 & lognormal & C & no \\
nuc-G2 & E & 57 & 5 & lognormal & C & no \\
nuc-G2 & L & 58 & 5 & lognormal & C & no \\
nuc-G3 & E & 59 & 5 & normal & C & yes \\
nuc-G4 & E & 61 & 5 & lognormal & C & no \\
nuc-G4 & L & 62 & 5 & lognormal & C & no \\
pdz-cis & EE & 1 & 5 & normal & NC & no \\
pdz-cis & ES & 4 & 5 & normal & NC & no \\
pdz-cis & LE & 3 & 5 & normal & NC & yes \\
pdz-cript-c & EE & 70 & 5 & normal & NC & no \\
pdz-cript-c & ES & 73 & 5 & normal & NC & no \\
pdz-cript-c & GE & 72 & 5 & lognormal & NC & no \\
pdz-cript-c & LE & 71 & 5 & normal & NC & no \\
pdz-cript-c & LS & 74 & 5 & normal & NC & no \\
pdz-cript-n & EE & 76 & 5 & normal & NC & no \\
pdz-cript-n & GE & 78 & 5 & lognormal & NC & no \\
pdz-cript-n & GS & 81 & 5 & lognormal & NC & no \\
pdz-cript-n & LE & 77 & 5 & normal & NC & no \\
pdz-cript-n & LS & 80 & 5 & normal & NC & no \\
pdz-trans-1 & EE & 82 & 5 & normal & NC & no \\
pdz-trans-1 & GE & 84 & 5 & lognormal & NC & no \\
pdz-trans-2 & EE & 88 & 5 & normal & NC & no \\
pdz-trans-2 & GE & 90 & 5 & lognormal & NC & no \\
pdz-trans-2 & GS & 93 & 5 & lognormal & NC & no \\
pdz-trans-2 & LS & 92 & 5 & normal & NC & no \\
trpb-four-site & EE & 64 & 5 & normal & NC & yes \\
trpb-four-site & ES & 67 & 5 & normal & NC & no \\
trpb-four-site & GE & 66 & 5 & lognormal & NC & no \\
trpb-four-site & GS & 69 & 5 & lognormal & NC & no \\
trpb-four-site & LE & 65 & 5 & normal & NC & no \\
trpb-four-site & LS & 68 & 5 & normal & NC & no \\
trpb-libA & EE & 7 & 0.1 & lognormal & NC & no \\
trpb-libA & ES & 9 & 0.1 & lognormal & NC & no \\
trpb-libA & LE & 8 & 0.1 & lognormal & NC & no \\
trpb-libB & EE & 11 & 0.1 & lognormal & NC & no \\
trpb-libB & ES & 13 & 0.1 & lognormal & NC & no \\
trpb-libB & LE & 12 & 0.1 & lognormal & NC & no \\
trpb-libC & EE & 15 & 0.1 & lognormal & NC & no \\
trpb-libC & ES & 17 & 0.1 & lognormal & NC & no \\
trpb-libC & LE & 16 & 0.1 & lognormal & NC & no \\
trpb-libD & EE & 19 & 0.1 & normal & NC & no \\
trpb-libD & GE & 21 & 0.1 & lognormal & NC & no \\
trpb-libD & LE & 20 & 0.1 & normal & NC & no \\
trpb-libE & EE & 25 & 0.1 & lognormal & C & no \\
trpb-libE & GE & 27 & 0.1 & lognormal & NC & no \\
trpb-libE & LE & 26 & 0.1 & normal & NC & no \\
trpb-libF & EE & 31 & 0.1 & normal & NC & no \\
trpb-libF & ES & 34 & 0.1 & normal & NC & no \\
trpb-libF & GE & 33 & 0.1 & lognormal & NC & no \\
trpb-libF & LE & 32 & 0.1 & normal & NC & no \\
trpb-libG & EE & 37 & 0.1 & normal & NC & no \\
trpb-libG & GE & 39 & 0.1 & lognormal & NC & no \\
trpb-libG & LE & 38 & 0.1 & normal & NC & no \\
trpb-libH & EE & 43 & 0.1 & normal & NC & no \\
trpb-libH & GE & 45 & 0.1 & lognormal & NC & no \\
trpb-libH & LE & 44 & 0.1 & normal & NC & no \\
trpb-libI & EE & 49 & 0.1 & normal & NC & no \\
trpb-libI & GE & 51 & 0.1 & lognormal & NC & no \\
trpb-libI & LE & 50 & 0.1 & normal & NC & no \\
\end{longtable}
\endgroup


\note{TODO: Supply versioned public archives and complete environment locks,
and assemble the remaining fit-specific predictive and sensitivity exhibits.
Full-landscape uncertainty estimation, read-depth interventions, and FLIP2
predictor evaluation remain outside this draft.}
